\chapter{Nota Introduttiva}

Poco si conosce dell'origine della Dottrina del Fiume e del libro che la espone, conosciuto come
Il Libro dei Tre Anelli. La dottrina in sé non fa parte delle scuole di filosofia conosciute,
e non sono noti importanti e rappresentativi studiosi, nemmeno accademie, che l'hanno adottata
come propria.
Nel testo originale è presente una nota dei presunti fondatori che sottolinea
la volontà di questi di diffonderne il pensiero e la pratica da essa derivata, detta  "fluminautica", senza clamore o proselitismi.\par

Sono state proposte numerose spiegazioni sulla sua origine, la più nota è quella che
vuole alcuni discepoli di Sharess allontanarsi dal culto dei festali, insoddisfatti dalle
risposte vaghe e caotiche. Si pensa che la dottrina sia scaturita da una riflessione sul ruolo
che le emozioni ricoprono nella vita di un individuo, ma rimane una speculazione che non ha
trovato conferme o smentite.\par

Questa edizione raccoglie il testo conservato a Candlekeep, aggiungendo note di ricerca
filologica e storica, conquistate con fatica nel corso degli anni, di questa dottrina semisconosciuta.
Va notato che, non esistendo documenti specifici, tale ricerca ha dovuto ricorrere a metodi volti a
ricostruire piuttosto che registrare, per tanto si sono cercate tracce della dottrina nelle pratiche giuridiche,
nel tipo di riforme sociali, e nei racconti e nelle poche testimonianze riportate di seconda mano
da avventurieri e viaggiatori generosi.

Va da sé che la natura di questa ricostruzione rimane speculativa.
Non riteniamo che sia un difetto, specialmente se vogliamo rispettare il volere dei suoi autori, che aspiravano
a una diffusione gentile e meritocratica della loro creatura.
\cleardoublepage